\documentclass[12pt,a4paper]{article}
\usepackage[utf8]{inputenc}
\usepackage[spanish]{babel}
\usepackage{geometry}
\usepackage{longtable}
\usepackage{booktabs}
\usepackage{xcolor}
\usepackage{hyperref}
\usepackage{array}
\usepackage{multirow}
\usepackage{enumitem}

\geometry{margin=2.5cm}
\hypersetup{colorlinks=true, linkcolor=blue, urlcolor=blue}

\definecolor{navy}{HTML}{1B2A4A}
\definecolor{teal}{HTML}{0D9488}

% --- Comando para ficha de tabla ---
\newenvironment{tabladesc}[3]{%
  \section{#1}
  \textbf{Nombre físico:} \texttt{#2} \hfill \textbf{Tipo:} #3\par
  \medskip
}{%
  \bigskip
}

\newcolumntype{L}[1]{>{\raggedright\arraybackslash}p{#1}}

\begin{document}

\begin{titlepage}
\centering
\vspace*{4cm}
{\Huge\color{navy}\textbf{Diccionario de Datos}}\\[0.8cm]
{\Large Sistema de Gestión de Laboratorios --- SIGMALAB}\\[0.4cm]
{\large Universidad Mayor de San Andrés}\\[2cm]
\vfill
{\small Generado desde \texttt{schema.prisma} --- \today}
\end{titlepage}

\tableofcontents
\newpage

%============================================================
\section{Introducción}
%============================================================

Este documento describe la estructura completa de la base de datos del sistema SIGMALAB.
Se definen 24 tablas que cubren los módulos de: inventario de equipos, periféricos,
mantenimientos (preventivo y correctivo), incidencias, reportes de pasantes,
asignaciones, usuarios, roles, auditoría y trazabilidad.

Convenciones utilizadas:
\begin{itemize}[nosep]
  \item \texttt{PK} --- Clave primaria.
  \item \texttt{FK} --- Clave foránea.
  \item \texttt{UQ} --- Restricción de unicidad.
  \item \texttt{NN} --- No nulo (\texttt{NOT NULL}).
  \item \texttt{DF} --- Valor por defecto.
\end{itemize}

\newpage

%============================================================
% TABLA 1
%============================================================
\begin{tabladesc}{Edificio}{edificios}{Maestra}

Catalogación de edificios que albergan los laboratorios.

\bigskip
\begin{longtable}{L{3.5cm} L{2cm} L{1.8cm} L{8cm}}
\toprule
\textbf{Atributo} & \textbf{Tipo} & \textbf{Restricciones} & \textbf{Descripción} \\
\midrule
\endhead
\bottomrule
\endfoot
\texttt{id}       & VARCHAR(10)  & PK, NN & Identificador único del edificio (ej. \texttt{ED-001}). \\
\texttt{nombre}   & VARCHAR(100) & NN     & Nombre completo del edificio. \\
\texttt{ubicacion} & VARCHAR(200)& ---    & Dirección o referencia de ubicación. \\
\texttt{activo}   & BOOLEAN      & DF: true & Indica si el registro está activo. \\
\bottomrule
\end{longtable}

\textbf{Relaciones:}
\begin{itemize}[nosep]
  \item 1 $\rightarrow$ N con \texttt{Laboratorio} (\texttt{laboratorios}).
\end{itemize}
\end{tabladesc}

%============================================================
% TABLA 2
%============================================================
\begin{tabladesc}{Rol}{roles}{Maestra}

Catálogo de roles del sistema para control de acceso.

\bigskip
\begin{longtable}{L{3.5cm} L{2cm} L{1.8cm} L{8cm}}
\toprule
\textbf{Atributo} & \textbf{Tipo} & \textbf{Restricciones} & \textbf{Descripción} \\
\midrule
\endhead
\bottomrule
\endfoot
\texttt{id}         & VARCHAR(20)  & PK, NN & Identificador del rol (ej. \texttt{encargado}). \\
\texttt{nombre}     & VARCHAR(50)  & UQ, NN & Nombre legible del rol. \\
\texttt{descripcion}& VARCHAR(200) & ---    & Descripción del alcance del rol. \\
\texttt{nivelAcceso}& SMALLINT     & NN     & Nivel jerárquico (mayor valor = más privilegios). \\
\bottomrule
\end{longtable}

\textbf{Relaciones:}
\begin{itemize}[nosep]
  \item 1 $\rightarrow$ N con \texttt{Usuario} (\texttt{usuarios}).
\end{itemize}
\end{tabladesc}

%============================================================
% TABLA 3
%============================================================
\begin{tabladesc}{Categoría de Inventario}{categorias\_inventario}{Maestra}

Categorías para clasificar los ítems de inventario (consumibles, repuestos, etc.).

\bigskip
\begin{longtable}{L{3.5cm} L{2cm} L{1.8cm} L{8cm}}
\toprule
\textbf{Atributo} & \textbf{Tipo} & \textbf{Restricciones} & \textbf{Descripción} \\
\midrule
\endhead
\bottomrule
\endfoot
\texttt{id}        & VARCHAR(30)  & PK, NN & Identificador único. \\
\texttt{nombre}    & VARCHAR(100) & UQ, NN & Nombre de la categoría. \\
\texttt{stockMinimo}& INT         & DF: 0  & Umbral mínimo de stock para alertas. \\
\texttt{activo}    & BOOLEAN      & DF: true & Indica si la categoría está vigente. \\
\bottomrule
\end{longtable}

\textbf{Relaciones:}
\begin{itemize}[nosep]
  \item 1 $\rightarrow$ N con \texttt{InventarioItem} (\texttt{inventario}).
\end{itemize}
\end{tabladesc}

%============================================================
% TABLA 4
%============================================================
\begin{tabladesc}{Tipo de Mantenimiento}{tipo\_mantenimiento}{Maestra}

Clasificación del tipo de mantenimiento (preventivo, correctivo, predictivo).

\bigskip
\begin{longtable}{L{3.5cm} L{2cm} L{1.8cm} L{8cm}}
\toprule
\textbf{Atributo} & \textbf{Tipo} & \textbf{Restricciones} & \textbf{Descripción} \\
\midrule
\endhead
\bottomrule
\endfoot
\texttt{id}     & VARCHAR(20) & PK, NN & Identificador del tipo. \\
\texttt{nombre} & VARCHAR(50) & UQ, NN & Nombre del tipo de mantenimiento. \\
\bottomrule
\end{longtable}

\textbf{Relaciones:}
\begin{itemize}[nosep]
  \item 1 $\rightarrow$ N con \texttt{Mantenimiento} (\texttt{mantenimientos}).
\end{itemize}
\end{tabladesc}

%============================================================
% TABLA 5
%============================================================
\begin{tabladesc}{Estado de Mantenimiento}{estado\_mantenimiento}{Maestra}

Estados posibles del flujo de un mantenimiento.

\bigskip
\begin{longtable}{L{3.5cm} L{2cm} L{1.8cm} L{8cm}}
\toprule
\textbf{Atributo} & \textbf{Tipo} & \textbf{Restricciones} & \textbf{Descripción} \\
\midrule
\endhead
\bottomrule
\endfoot
\texttt{id}     & VARCHAR(20) & PK, NN & Identificador del estado. \\
\texttt{nombre} & VARCHAR(50) & UQ, NN & Nombre del estado. \\
\bottomrule
\end{longtable}

\textbf{Relaciones:}
\begin{itemize}[nosep]
  \item 1 $\rightarrow$ N con \texttt{Mantenimiento} (\texttt{mantenimientos}).
\end{itemize}
\end{tabladesc}

%============================================================
% TABLA 6
%============================================================
\begin{tabladesc}{Estado de Equipo}{estado\_equipo}{Maestra}

Catálogo de estados operativos de un equipo (Operativo, En mantenimiento, De baja, etc.).

\bigskip
\begin{longtable}{L{3.5cm} L{2cm} L{1.8cm} L{8cm}}
\toprule
\textbf{Atributo} & \textbf{Tipo} & \textbf{Restricciones} & \textbf{Descripción} \\
\midrule
\endhead
\bottomrule
\endfoot
\texttt{id}     & VARCHAR(30) & PK, NN & Identificador del estado. \\
\texttt{nombre} & VARCHAR(50) & UQ, NN & Nombre del estado. \\
\bottomrule
\end{longtable}

\textbf{Relaciones:}
\begin{itemize}[nosep]
  \item 1 $\rightarrow$ N con \texttt{Equipo} (\texttt{equipos}).
\end{itemize}
\end{tabladesc}

%============================================================
% TABLA 7
%============================================================
\begin{tabladesc}{Estado de Incidencia}{estado\_incidencia}{Maestra}

Estados del ciclo de vida de una incidencia (Nuevo, En proceso, Resuelto, etc.).

\bigskip
\begin{longtable}{L{3.5cm} L{2cm} L{1.8cm} L{8cm}}
\toprule
\textbf{Atributo} & \textbf{Tipo} & \textbf{Restricciones} & \textbf{Descripción} \\
\midrule
\endhead
\bottomrule
\endfoot
\texttt{id}     & VARCHAR(20) & PK, NN & Identificador del estado. \\
\texttt{nombre} & VARCHAR(50) & UQ, NN & Nombre del estado. \\
\bottomrule
\end{longtable}

\textbf{Relaciones:}
\begin{itemize}[nosep]
  \item 1 $\rightarrow$ N con \texttt{Incidencia} (\texttt{incidencias}).
\end{itemize}
\end{tabladesc}

%============================================================
% TABLA 8
%============================================================
\begin{tabladesc}{Persona}{personas}{Maestra}

Datos personales de individuos vinculados al sistema (docentes, estudiantes, técnicos, encargados).

\bigskip
\begin{longtable}{L{3.5cm} L{2.5cm} L{2cm} L{7.5cm}}
\toprule
\textbf{Atributo} & \textbf{Tipo} & \textbf{Restricciones} & \textbf{Descripción} \\
\midrule
\endhead
\bottomrule
\endfoot
\texttt{id}                & VARCHAR(20)  & PK, NN & Identificador único. \\
\texttt{nombres}           & VARCHAR(100) & NN     & Nombres de la persona. \\
\texttt{paterno}           & VARCHAR(50)  & NN     & Apellido paterno. \\
\texttt{materno}           & VARCHAR(50)  & ---    & Apellido materno. \\
\texttt{ci}                & VARCHAR(20)  & UQ     & Cédula de identidad. \\
\texttt{direccion}         & VARCHAR(255) & ---    & Dirección de domicilio. \\
\texttt{registroUniversitario} & VARCHAR(20) & UQ  & Registro universitario (RU). \\
\texttt{email}             & VARCHAR(150) & UQ     & Correo electrónico. \\
\texttt{celular}           & VARCHAR(20)  & ---    & Número de teléfono celular. \\
\texttt{fotoUrl}           & TEXT         & ---    & URL de la foto de perfil. \\
\texttt{fechaIngresoPasante}& DATE        & ---    & Fecha de ingreso del pasante. \\
\texttt{activo}            & BOOLEAN      & DF: true & Registro activo. \\
\texttt{createdAt}         & TIMESTAMP    & NN, DF: now & Fecha de creación. \\
\texttt{updatedAt}         & TIMESTAMP    & NN, DF: now & Fecha de última modificación. \\
\bottomrule
\end{longtable}

\textbf{Relaciones:}
\begin{itemize}[nosep]
  \item 1 $\leftrightarrow$ 1 con \texttt{Usuario} (\texttt{usuario}) — Cascade en delete.
  \item 1 $\rightarrow$ N con \texttt{Laboratorio} (\texttt{laboratorios}) como encargado.
  \item 1 $\rightarrow$ N con \texttt{Incidencia} (\texttt{incidencias}) como persona reportante.
\end{itemize}
\end{tabladesc}

%============================================================
% TABLA 9
%============================================================
\begin{tabladesc}{Usuario}{usuarios}{Maestra / Operacional}

Cuentas de acceso al sistema. Cada usuario tiene una persona asociada.

\bigskip
\begin{longtable}{L{3.5cm} L{2.5cm} L{2cm} L{7.5cm}}
\toprule
\textbf{Atributo} & \textbf{Tipo} & \textbf{Restricciones} & \textbf{Descripción} \\
\midrule
\endhead
\bottomrule
\endfoot
\texttt{id}          & VARCHAR(20)  & PK, NN & Identificador único (username). \\
\texttt{personaId}   & VARCHAR(20)  & FK, UQ, NN & Referencia a \texttt{Persona.id}. \\
\texttt{roleId}      & VARCHAR(20)  & FK, NN & Referencia a \texttt{Rol.id}. \\
\texttt{passwordHash}& VARCHAR(255) & NN     & Hash bcrypt de la contraseña. \\
\texttt{ultimoAcceso}& TIMESTAMP    & ---    & Último inicio de sesión. \\
\texttt{activo}      & BOOLEAN      & DF: true & Cuenta habilitada. \\
\texttt{createdAt}   & TIMESTAMP    & NN, DF: now & Fecha de creación. \\
\texttt{updatedAt}   & TIMESTAMP    & NN, DF: now & Fecha de última modificación. \\
\bottomrule
\end{longtable}

\textbf{Relaciones:}
\begin{itemize}[nosep]
  \item 1 $\leftrightarrow$ 1 con \texttt{Persona} (\texttt{persona}).
  \item N $\rightarrow$ 1 con \texttt{Rol} (\texttt{rol}).
  \item 1 $\rightarrow$ N con \texttt{Asignacion} (\texttt{asignaciones}).
  \item 1 $\rightarrow$ N con \texttt{Incidencia} (\texttt{incidencias}).
  \item 1 $\rightarrow$ N con \texttt{Log} (\texttt{logs}).
  \item 1 $\rightarrow$ N con \texttt{AuditChange} (\texttt{auditChanges}).
  \item 1 $\rightarrow$ N con \texttt{Mantenimiento} (\texttt{mantenimientos}).
  \item 1 $\rightarrow$ N con \texttt{ReportePasante} como atendido por (\texttt{reportesAtendidos}).
  \item 1 $\rightarrow$ N con \texttt{ReportePasante} como pasante (\texttt{reportesPasante}).
  \item 1 $\rightarrow$ N con \texttt{PasanteHistorial} (\texttt{pasanteHistorial}).
\end{itemize}
\end{tabladesc}

%============================================================
% TABLA 10
%============================================================
\begin{tabladesc}{Historial de Pasante}{pasante\_historial}{Histórica}

Registro de cambios en la fecha de ingreso de un pasante.

\bigskip
\begin{longtable}{L{3.5cm} L{2.5cm} L{2cm} L{7.5cm}}
\toprule
\textbf{Atributo} & \textbf{Tipo} & \textbf{Restricciones} & \textbf{Descripción} \\
\midrule
\endhead
\bottomrule
\endfoot
\texttt{id}          & VARCHAR(30) & PK, NN & Identificador único. \\
\texttt{usuarioId}   & VARCHAR(20) & FK, NN & Referencia a \texttt{Usuario.id}. \\
\texttt{fechaIngreso}& DATE        & NN     & Nueva fecha de ingreso registrada. \\
\texttt{fechaCambio} & TIMESTAMP   & NN, DF: now & Momento del cambio. \\
\texttt{modificadoPor}& VARCHAR(20)& ---    & ID del usuario que realizó el cambio. \\
\texttt{motivo}      & VARCHAR(200)& ---    & Razón del cambio. \\
\bottomrule
\end{longtable}

\textbf{Relaciones:}
\begin{itemize}[nosep]
  \item N $\rightarrow$ 1 con \texttt{Usuario} (\texttt{usuario}).
\end{itemize}
\end{tabladesc}

%============================================================
% TABLA 11
%============================================================
\begin{tabladesc}{Laboratorio}{laboratorios}{Maestra / Operacional}

Espacios físicos (laboratorios) donde se ubican los equipos.

\bigskip
\begin{longtable}{L{3.5cm} L{2.5cm} L{2cm} L{7.5cm}}
\toprule
\textbf{Atributo} & \textbf{Tipo} & \textbf{Restricciones} & \textbf{Descripción} \\
\midrule
\endhead
\bottomrule
\endfoot
\texttt{id}              & VARCHAR(20)  & PK, NN & Identificador único (ej. \texttt{LAB-001}). \\
\texttt{nombre}          & VARCHAR(100) & NN     & Nombre del laboratorio. \\
\texttt{edificioId}      & VARCHAR(10)  & FK, NN & Referencia a \texttt{Edificio.id}. \\
\texttt{piso}            & SMALLINT     & NN     & Número de piso. \\
\texttt{capacidadEquipos}& INT          & NN     & Cantidad máxima de equipos. \\
\texttt{capacidadPersonas}& INT         & NN     & Capacidad de personas. \\
\texttt{encargadoId}     & VARCHAR(20)  & FK     & Referencia a \texttt{Persona.id} (encargado). \\
\texttt{activo}          & BOOLEAN      & DF: true & Laboratorio habilitado. \\
\texttt{createdAt}       & TIMESTAMP    & NN, DF: now & Fecha de creación. \\
\texttt{updatedAt}       & TIMESTAMP    & NN, DF: now & Fecha de última modificación. \\
\bottomrule
\end{longtable}

\textbf{Relaciones:}
\begin{itemize}[nosep]
  \item N $\rightarrow$ 1 con \texttt{Edificio} (\texttt{edificio}).
  \item N $\rightarrow$ 1 con \texttt{Persona} (\texttt{encargado}).
  \item 1 $\rightarrow$ N con \texttt{Equipo} (\texttt{equipos}).
  \item 1 $\rightarrow$ N con \texttt{Periferico} (\texttt{perifericos}).
  \item 1 $\rightarrow$ N con \texttt{InventarioItem} (\texttt{inventario}).
  \item 1 $\rightarrow$ N con \texttt{Incidencia} (\texttt{incidencias}).
  \item 1 $\rightarrow$ N con \texttt{Mantenimiento} (\texttt{mantenimientos}).
  \item 1 $\rightarrow$ N con \texttt{Asignacion} (\texttt{asignaciones}).
  \item 1 $\rightarrow$ N con \texttt{ReportePasante} (\texttt{reportesPasante}).
\end{itemize}
\end{tabladesc}

%============================================================
% TABLA 12
%============================================================
\begin{tabladesc}{Equipo}{equipos}{Maestra / Operacional}

Equipos de cómputo instalados en los laboratorios.

\bigskip
\begin{longtable}{L{3.5cm} L{2.5cm} L{2cm} L{7.5cm}}
\toprule
\textbf{Atributo} & \textbf{Tipo} & \textbf{Restricciones} & \textbf{Descripción} \\
\midrule
\endhead
\bottomrule
\endfoot
\texttt{codigo}          & VARCHAR(30)  & PK, NN & Código único del equipo (ej. \texttt{EQ-001}). \\
\texttt{nombre}          & VARCHAR(150) & NN     & Nombre o descripción del equipo. \\
\texttt{laboratorioId}   & VARCHAR(20)  & FK, NN & Referencia a \texttt{Laboratorio.id}. \\
\texttt{fila}            & VARCHAR(5)   & ---    & Fila dentro del laboratorio. \\
\texttt{puesto}          & VARCHAR(5)   & ---    & Número de puesto. \\
\texttt{sistemaOperativo}& VARCHAR(100) & ---    & SO instalado. \\
\texttt{marca}           & VARCHAR(50)  & ---    & Marca del equipo. \\
\texttt{modelo}          & VARCHAR(100) & ---    & Modelo del equipo. \\
\texttt{numeroSerie}     & VARCHAR(100) & UQ     & Número de serie del fabricante. \\
\texttt{codigoItic}      & VARCHAR(50)  & UQ     & Código de inventario ITIC. \\
\texttt{codigoFacultativo}& VARCHAR(50) & ---    & Código de inventario facultativo. \\
\texttt{codigoUmsa}      & VARCHAR(50)  & ---    & Código de inventario UMSA. \\
\texttt{estadoId}        & VARCHAR(30)  & FK, NN & Referencia a \texttt{EstadoEquipo.id}. \\
\texttt{fechaCompra}     & DATE         & ---    & Fecha de adquisición. \\
\texttt{fechaBaja}       & DATE         & ---    & Fecha de baja del equipo. \\
\texttt{motivoBaja}      & TEXT         & ---    & Razón de la baja. \\
\texttt{reemplazadoPor}  & VARCHAR(30)  & ---    & Código del equipo que lo reemplazó. \\
\texttt{activo}          & BOOLEAN      & DF: true & Equipo habilitado. \\
\texttt{createdAt}       & TIMESTAMP    & NN, DF: now & Fecha de creación. \\
\texttt{updatedAt}       & TIMESTAMP    & NN, DF: now & Fecha de última modificación. \\
\midrule
\multicolumn{4}{l}{\small\textit{Restricción única compuesta:}} \\
\multicolumn{4}{l}{\small\texttt{UQ(laboratorioId, fila, puesto)}} \\
\bottomrule
\end{longtable}

\textbf{Relaciones:}
\begin{itemize}[nosep]
  \item N $\rightarrow$ 1 con \texttt{EstadoEquipo} (\texttt{estado}).
  \item N $\rightarrow$ 1 con \texttt{Laboratorio} (\texttt{laboratorio}).
  \item 1 $\rightarrow$ N con \texttt{Asignacion} (\texttt{asignaciones}).
  \item 1 $\rightarrow$ N con \texttt{Incidencia} (\texttt{incidencias}).
  \item 1 $\rightarrow$ N con \texttt{InventarioItem} (\texttt{inventarioItems}).
  \item 1 $\rightarrow$ N con \texttt{Mantenimiento} (\texttt{mantenimientos}).
  \item 1 $\rightarrow$ N con \texttt{Periferico} (\texttt{perifericos}).
\end{itemize}
\end{tabladesc}

%============================================================
% TABLA 13
%============================================================
\begin{tabladesc}{Periférico}{perifericos}{Operacional}

Dispositivos periféricos (monitores, teclados, mouse, etc.) asociados a equipos o laboratorios.

\bigskip
\begin{longtable}{L{3.5cm} L{2.5cm} L{2cm} L{7.5cm}}
\toprule
\textbf{Atributo} & \textbf{Tipo} & \textbf{Restricciones} & \textbf{Descripción} \\
\midrule
\endhead
\bottomrule
\endfoot
\texttt{id}              & VARCHAR(30)  & PK, NN & Identificador único. \\
\texttt{tipo}            & VARCHAR(50)  & NN     & Tipo de periférico (Monitor, Teclado, etc.). \\
\texttt{marca}           & VARCHAR(50)  & ---    & Marca del periférico. \\
\texttt{modelo}          & VARCHAR(100) & ---    & Modelo del periférico. \\
\texttt{numeroSerie}     & VARCHAR(100) & ---    & Número de serie. \\
\texttt{codigoItic}      & VARCHAR(50)  & UQ     & Código de inventario ITIC. \\
\texttt{codigoFacultativo}& VARCHAR(50) & ---    & Código de inventario facultativo. \\
\texttt{codigoUmsa}      & VARCHAR(50)  & ---    & Código de inventario UMSA. \\
\texttt{laboratorioId}   & VARCHAR(20)  & FK     & Referencia a \texttt{Laboratorio.id}. \\
\texttt{equipoCodigo}    & VARCHAR(30)  & FK     & Referencia a \texttt{Equipo.codigo}. \\
\texttt{estado}          & VARCHAR(30)  & NN, DF: 'Funcionando' & Estado operativo. \\
\texttt{activo}          & BOOLEAN      & DF: true & Registro activo. \\
\texttt{createdAt}       & TIMESTAMP    & NN, DF: now & Fecha de creación. \\
\bottomrule
\end{longtable}

\textbf{Relaciones:}
\begin{itemize}[nosep]
  \item N $\rightarrow$ 1 con \texttt{Equipo} (\texttt{equipo}).
  \item N $\rightarrow$ 1 con \texttt{Laboratorio} (\texttt{laboratorio}).
\end{itemize}
\end{tabladesc}

%============================================================
% TABLA 14
%============================================================
\begin{tabladesc}{Ítem de Inventario}{inventario}{Operacional}

Control de inventario de consumibles, repuestos y materiales.

\bigskip
\begin{longtable}{L{3.5cm} L{2.5cm} L{2cm} L{7.5cm}}
\toprule
\textbf{Atributo} & \textbf{Tipo} & \textbf{Restricciones} & \textbf{Descripción} \\
\midrule
\endhead
\bottomrule
\endfoot
\texttt{id}              & VARCHAR(30)  & PK, NN & Identificador único. \\
\texttt{categoriaId}     & VARCHAR(30)  & FK, NN & Referencia a \texttt{CategoriaInventario.id}. \\
\texttt{codigoItic}      & VARCHAR(50)  & UQ, NN & Código de inventario ITIC. \\
\texttt{codigoFacultativo}& VARCHAR(50) & ---    & Código de inventario facultativo. \\
\texttt{codigoUmsa}      & VARCHAR(50)  & ---    & Código de inventario UMSA. \\
\texttt{numeroSerie}     & VARCHAR(100) & ---    & Número de serie. \\
\texttt{marca}           & VARCHAR(50)  & ---    & Marca del ítem. \\
\texttt{modelo}          & VARCHAR(100) & ---    & Modelo del ítem. \\
\texttt{estado}          & VARCHAR(30)  & NN, DF: 'En almacén' & Estado actual. \\
\texttt{fechaIngreso}    & DATE         & NN     & Fecha de ingreso al inventario. \\
\texttt{fechaAsignacion} & DATE         & ---    & Fecha de asignación a un equipo/lab. \\
\texttt{laboratorioId}   & VARCHAR(20)  & FK     & Referencia a \texttt{Laboratorio.id}. \\
\texttt{equipoCodigo}    & VARCHAR(30)  & FK     & Referencia a \texttt{Equipo.codigo}. \\
\texttt{observaciones}   & TEXT         & ---    & Notas adicionales. \\
\texttt{activo}          & BOOLEAN      & DF: true & Registro activo. \\
\texttt{createdAt}       & TIMESTAMP    & NN, DF: now & Fecha de creación. \\
\bottomrule
\end{longtable}

\textbf{Relaciones:}
\begin{itemize}[nosep]
  \item N $\rightarrow$ 1 con \texttt{CategoriaInventario} (\texttt{categoria}).
  \item N $\rightarrow$ 1 con \texttt{Equipo} (\texttt{equipo}).
  \item N $\rightarrow$ 1 con \texttt{Laboratorio} (\texttt{laboratorio}).
\end{itemize}
\end{tabladesc}

%============================================================
% TABLA 15
%============================================================
\begin{tabladesc}{Insumo}{insumos}{Maestra / Operacional}

Catálogo de insumos utilizados en mantenimientos.

\bigskip
\begin{longtable}{L{3.5cm} L{2.5cm} L{2cm} L{7.5cm}}
\toprule
\textbf{Atributo} & \textbf{Tipo} & \textbf{Restricciones} & \textbf{Descripción} \\
\midrule
\endhead
\bottomrule
\endfoot
\texttt{nombre}     & VARCHAR(100) & PK, NN & Nombre del insumo. \\
\texttt{tipo}       & VARCHAR(50)  & ---    & Clasificación del insumo. \\
\texttt{unidadMedida}& VARCHAR(30) & NN     & Unidad de medida (ej. \texttt{unidad}, \texttt{ml}). \\
\texttt{stock}      & INT          & DF: 0  & Cantidad disponible. \\
\texttt{stockMinimo}& INT          & DF: 0  & Umbral mínimo de stock. \\
\texttt{activo}     & BOOLEAN      & DF: true & Insumo habilitado. \\
\texttt{createdAt}  & TIMESTAMP    & NN, DF: now & Fecha de creación. \\
\texttt{updatedAt}  & TIMESTAMP    & NN, DF: now & Fecha de última modificación. \\
\bottomrule
\end{longtable}

\textbf{Relaciones:}
\begin{itemize}[nosep]
  \item 1 $\rightarrow$ N con \texttt{InsumoUsado} (\texttt{insumosUsados}).
\end{itemize}
\end{tabladesc}

%============================================================
% TABLA 16
%============================================================
\begin{tabladesc}{Mantenimiento}{mantenimientos}{Operacional}

Registro de mantenimientos realizados a equipos.

\bigskip
\begin{longtable}{L{3.5cm} L{2.5cm} L{2cm} L{7.5cm}}
\toprule
\textbf{Atributo} & \textbf{Tipo} & \textbf{Restricciones} & \textbf{Descripción} \\
\midrule
\endhead
\bottomrule
\endfoot
\texttt{id}          & VARCHAR(30)  & PK, NN & Identificador único. \\
\texttt{tipoId}      & VARCHAR(20)  & FK, NN & Referencia a \texttt{TipoMantenimiento.id}. \\
\texttt{equipoCodigo}& VARCHAR(30)  & FK, NN & Referencia a \texttt{Equipo.codigo}. \\
\texttt{tecnicoId}   & VARCHAR(20)  & FK, NN & Referencia a \texttt{Usuario.id}. \\
\texttt{laboratorioId}& VARCHAR(20) & FK     & Referencia a \texttt{Laboratorio.id}. \\
\texttt{fecha}       & DATE         & NN     & Fecha del mantenimiento. \\
\texttt{horaInicio}  & TIME         & ---    & Hora de inicio. \\
\texttt{horaFin}     & TIME         & ---    & Hora de finalización. \\
\texttt{estadoId}    & VARCHAR(20)  & FK, NN & Referencia a \texttt{EstadoMantenimiento.id}. \\
\texttt{activo}      & BOOLEAN      & DF: true & Registro activo. \\
\texttt{createdAt}   & TIMESTAMP    & NN, DF: now & Fecha de creación. \\
\texttt{updatedAt}   & TIMESTAMP    & NN, DF: now & Fecha de última modificación. \\
\bottomrule
\end{longtable}

\textbf{Relaciones:}
\begin{itemize}[nosep]
  \item N $\rightarrow$ 1 con \texttt{TipoMantenimiento} (\texttt{tipo}).
  \item N $\rightarrow$ 1 con \texttt{EstadoMantenimiento} (\texttt{estado}).
  \item N $\rightarrow$ 1 con \texttt{Equipo} (\texttt{equipo}).
  \item N $\rightarrow$ 1 con \texttt{Laboratorio} (\texttt{laboratorio}).
  \item N $\rightarrow$ 1 con \texttt{Usuario} (\texttt{tecnico}).
  \item 1 $\leftrightarrow$ 1 con \texttt{MantenimientoDetalle} (\texttt{detalle}).
\end{itemize}
\end{tabladesc}

%============================================================
% TABLA 17
%============================================================
\begin{tabladesc}{Detalle de Mantenimiento}{mantenimientos\_detalle}{Operacional}

Información detallada del diagnóstico, acciones y resoluciones de un mantenimiento.

\bigskip
\begin{longtable}{L{3.5cm} L{2.5cm} L{2cm} L{7.5cm}}
\toprule
\textbf{Atributo} & \textbf{Tipo} & \textbf{Restricciones} & \textbf{Descripción} \\
\midrule
\endhead
\bottomrule
\endfoot
\texttt{id}             & VARCHAR(30) & PK, NN & Identificador único. \\
\texttt{mantenimientoId}& VARCHAR(30) & FK, UQ, NN & Referencia a \texttt{Mantenimiento.id}. \\
\texttt{descripcion}    & TEXT        & ---    & Descripción general. \\
\texttt{diagnostico}    & TEXT        & ---    & Diagnóstico del problema. \\
\texttt{accionRealizada}& TEXT        & ---    & Acciones ejecutadas. \\
\texttt{resolucion}     & TEXT        & ---    & Resolución final. \\
\texttt{tipoIncidencia} & VARCHAR(50) & ---    & Tipo de incidencia detectada. \\
\texttt{estadoFinal}    & VARCHAR(50) & ---    & Estado del equipo después del mant. \\
\texttt{observaciones}  & TEXT        & ---    & Observaciones adicionales. \\
\texttt{recomendaciones}& TEXT        & ---    & Recomendaciones futuras. \\
\texttt{activo}         & BOOLEAN     & DF: true & Registro activo. \\
\bottomrule
\end{longtable}

\textbf{Relaciones:}
\begin{itemize}[nosep]
  \item 1 $\leftrightarrow$ 1 con \texttt{Mantenimiento} (\texttt{mantenimiento}) — Cascade en delete.
  \item 1 $\rightarrow$ N con \texttt{Checklist} (\texttt{checklists}) — Cascade en delete.
  \item 1 $\rightarrow$ N con \texttt{InsumoUsado} (\texttt{insumosUsados}) — Cascade en delete.
\end{itemize}
\end{tabladesc}

%============================================================
% TABLA 18
%============================================================
\begin{tabladesc}{Checklist}{checklists}{Operacional}

Items de verificación asociados a un mantenimiento.

\bigskip
\begin{longtable}{L{3.5cm} L{2.5cm} L{2cm} L{7.5cm}}
\toprule
\textbf{Atributo} & \textbf{Tipo} & \textbf{Restricciones} & \textbf{Descripción} \\
\midrule
\endhead
\bottomrule
\endfoot
\texttt{id}        & VARCHAR(30) & PK, NN & Identificador único. \\
\texttt{detalleId} & VARCHAR(30) & FK, NN & Referencia a \texttt{MantenimientoDetalle.id}. \\
\texttt{categoria} & VARCHAR(30) & NN     & Categoría del item (ej. \texttt{Hardware}). \\
\texttt{item}      & VARCHAR(200)& NN     & Descripción del elemento a verificar. \\
\texttt{estado}    & VARCHAR(20) & NN     & Estado (Bien, Mal, Regular, N/A). \\
\texttt{observacion}& TEXT       & ---    & Observación del técnico. \\
\texttt{activo}    & BOOLEAN     & DF: true & Registro activo. \\
\bottomrule
\end{longtable}

\textbf{Relaciones:}
\begin{itemize}[nosep]
  \item N $\rightarrow$ 1 con \texttt{MantenimientoDetalle} (\texttt{detalle}) — Cascade en delete.
\end{itemize}
\end{tabladesc}

%============================================================
% TABLA 19
%============================================================
\begin{tabladesc}{Insumo Usado}{insumos\_usados}{Operacional}

Registro de insumos consumidos en un mantenimiento.

\bigskip
\begin{longtable}{L{3.5cm} L{2.5cm} L{2cm} L{7.5cm}}
\toprule
\textbf{Atributo} & \textbf{Tipo} & \textbf{Restricciones} & \textbf{Descripción} \\
\midrule
\endhead
\bottomrule
\endfoot
\texttt{id}          & VARCHAR(30)  & PK, NN & Identificador único. \\
\texttt{detalleId}   & VARCHAR(30)  & FK, NN & Referencia a \texttt{MantenimientoDetalle.id}. \\
\texttt{insumoNombre}& VARCHAR(100) & FK, NN & Referencia a \texttt{Insumo.nombre}. \\
\texttt{cantidad}    & VARCHAR(50)  & NN     & Cantidad utilizada. \\
\texttt{activo}      & BOOLEAN      & DF: true & Registro activo. \\
\texttt{createdAt}   & TIMESTAMP    & NN, DF: now & Fecha de registro. \\
\bottomrule
\end{longtable}

\textbf{Relaciones:}
\begin{itemize}[nosep]
  \item N $\rightarrow$ 1 con \texttt{MantenimientoDetalle} (\texttt{detalle}) — Cascade en delete.
  \item N $\rightarrow$ 1 con \texttt{Insumo} (\texttt{insumo}).
\end{itemize}
\end{tabladesc}

%============================================================
% TABLA 20
%============================================================
\begin{tabladesc}{Incidencia}{incidencias}{Operacional}

Registro de incidentes reportados sobre equipos y laboratorios.

\bigskip
\begin{longtable}{L{3.5cm} L{2.5cm} L{2cm} L{7.5cm}}
\toprule
\textbf{Atributo} & \textbf{Tipo} & \textbf{Restricciones} & \textbf{Descripción} \\
\midrule
\endhead
\bottomrule
\endfoot
\texttt{id}                & VARCHAR(30)  & PK, NN & Identificador único. \\
\texttt{equipoCodigo}      & VARCHAR(30)  & FK     & Referencia a \texttt{Equipo.codigo}. \\
\texttt{laboratorioId}     & VARCHAR(20)  & FK     & Referencia a \texttt{Laboratorio.id}. \\
\texttt{usuarioId}         & VARCHAR(20)  & FK     & Referencia a \texttt{Usuario.id} (técnico). \\
\texttt{personaId}         & VARCHAR(20)  & FK     & Referencia a \texttt{Persona.id} (reportante). \\
\texttt{problema}          & TEXT         & NN     & Descripción del problema. \\
\texttt{requiereSeguimiento}& BOOLEAN     & DF: false & Indica si requiere seguimiento. \\
\texttt{estadoId}          & VARCHAR(20)  & FK, NN & Referencia a \texttt{EstadoIncidencia.id}. \\
\texttt{fecha}             & TIMESTAMP    & NN, DF: now & Fecha del reporte. \\
\texttt{resueltaEn}        & TIMESTAMP    & ---    & Fecha de resolución. \\
\texttt{activo}            & BOOLEAN      & DF: true & Registro activo. \\
\texttt{createdAt}         & TIMESTAMP    & NN, DF: now & Fecha de creación. \\
\texttt{updatedAt}         & TIMESTAMP    & NN, DF: now & Fecha de última modificación. \\
\bottomrule
\end{longtable}

\textbf{Relaciones:}
\begin{itemize}[nosep]
  \item N $\rightarrow$ 1 con \texttt{EstadoIncidencia} (\texttt{estado}).
  \item N $\rightarrow$ 1 con \texttt{Equipo} (\texttt{equipo}).
  \item N $\rightarrow$ 1 con \texttt{Laboratorio} (\texttt{laboratorio}).
  \item N $\rightarrow$ 1 con \texttt{Persona} (\texttt{persona}).
  \item N $\rightarrow$ 1 con \texttt{Usuario} (\texttt{usuario}).
\end{itemize}
\end{tabladesc}

%============================================================
% TABLA 21
%============================================================
\begin{tabladesc}{Asignación}{asignaciones}{Operacional}

Asignación de tareas técnicas a usuarios para atención de problemas.

\bigskip
\begin{longtable}{L{3.5cm} L{2.5cm} L{2cm} L{7.5cm}}
\toprule
\textbf{Atributo} & \textbf{Tipo} & \textbf{Restricciones} & \textbf{Descripción} \\
\midrule
\endhead
\bottomrule
\endfoot
\texttt{id}          & VARCHAR(30)  & PK, NN & Identificador único. \\
\texttt{equipoCodigo}& VARCHAR(30)  & FK, NN & Referencia a \texttt{Equipo.codigo}. \\
\texttt{laboratorioId}& VARCHAR(20) & FK     & Referencia a \texttt{Laboratorio.id}. \\
\texttt{tecnicoId}   & VARCHAR(20)  & FK, NN & Referencia a \texttt{Usuario.id}. \\
\texttt{problema}    & TEXT         & NN     & Descripción del problema asignado. \\
\texttt{prioridad}   & VARCHAR(10)  & NN, DF: 'Media' & Prioridad (Baja, Media, Alta). \\
\texttt{fecha}       & DATE         & NN, DF: now & Fecha de asignación. \\
\texttt{estado}      & VARCHAR(20)  & NN, DF: 'Pendiente' & Estado de la asignación. \\
\texttt{activo}      & BOOLEAN      & DF: true & Registro activo. \\
\texttt{createdAt}   & TIMESTAMP    & NN, DF: now & Fecha de creación. \\
\texttt{updatedAt}   & TIMESTAMP    & NN, DF: now & Fecha de última modificación. \\
\bottomrule
\end{longtable}

\textbf{Relaciones:}
\begin{itemize}[nosep]
  \item N $\rightarrow$ 1 con \texttt{Equipo} (\texttt{equipo}).
  \item N $\rightarrow$ 1 con \texttt{Laboratorio} (\texttt{laboratorio}).
  \item N $\rightarrow$ 1 con \texttt{Usuario} (\texttt{tecnico}).
\end{itemize}
\end{tabladesc}

%============================================================
% TABLA 22
%============================================================
\begin{tabladesc}{Reporte de Pasante}{reportes\_pasante}{Operacional}

Reportes generados por pasantes sobre novedades en laboratorios.

\bigskip
\begin{longtable}{L{3.5cm} L{2.5cm} L{2cm} L{7.5cm}}
\toprule
\textbf{Atributo} & \textbf{Tipo} & \textbf{Restricciones} & \textbf{Descripción} \\
\midrule
\endhead
\bottomrule
\endfoot
\texttt{id}             & VARCHAR(30)  & PK, NN & Identificador único. \\
\texttt{pasanteId}      & VARCHAR(20)  & FK, NN & Referencia a \texttt{Usuario.id} (pasante). \\
\texttt{rolReporte}     & VARCHAR(20)  & ---    & Rol del pasante al reportar. \\
\texttt{titulo}         & VARCHAR(200) & NN     & Título del reporte. \\
\texttt{descripcion}    & TEXT         & NN     & Descripción detallada. \\
\texttt{laboratorioId}  & VARCHAR(20)  & FK     & Referencia a \texttt{Laboratorio.id}. \\
\texttt{ubicacion}      & VARCHAR(200) & ---    & Ubicación específica. \\
\texttt{categoria}      & VARCHAR(50)  & ---    & Categoría del reporte. \\
\texttt{prioridad}      & VARCHAR(10)  & NN, DF: 'Media' & Prioridad. \\
\texttt{fecha}          & TIMESTAMP    & NN, DF: now & Fecha del reporte. \\
\texttt{estado}         & VARCHAR(30)  & NN, DF: 'Nuevo' & Estado del reporte. \\
\texttt{resolucionDetalle}& TEXT       & ---    & Detalle de la resolución. \\
\texttt{atendidoPor}    & VARCHAR(20)  & FK     & Referencia a \texttt{Usuario.id} (quien atendió). \\
\texttt{activo}         & BOOLEAN      & DF: true & Registro activo. \\
\texttt{createdAt}      & TIMESTAMP    & NN, DF: now & Fecha de creación. \\
\texttt{updatedAt}      & TIMESTAMP    & NN, DF: now & Fecha de última modificación. \\
\bottomrule
\end{longtable}

\textbf{Relaciones:}
\begin{itemize}[nosep]
  \item N $\rightarrow$ 1 con \texttt{Usuario} (\texttt{pasante}).
  \item N $\rightarrow$ 1 con \texttt{Usuario} (\texttt{atendidoPorUser}) como quien atiende.
  \item N $\rightarrow$ 1 con \texttt{Laboratorio} (\texttt{laboratorio}).
\end{itemize}
\end{tabladesc}

%============================================================
% TABLA 23
%============================================================
\begin{tabladesc}{Log}{logs}{Auditoría}

Registro de eventos y acciones del sistema para trazabilidad.

\bigskip
\begin{longtable}{L{3.5cm} L{2.5cm} L{2cm} L{7.5cm}}
\toprule
\textbf{Atributo} & \textbf{Tipo} & \textbf{Restricciones} & \textbf{Descripción} \\
\midrule
\endhead
\bottomrule
\endfoot
\texttt{id}          & VARCHAR(30)  & PK, NN & Identificador único. \\
\texttt{timestamp}   & TIMESTAMP    & NN, DF: now & Momento del evento. \\
\texttt{usuarioId}   & VARCHAR(20)  & FK     & Referencia a \texttt{Usuario.id}. \\
\texttt{accion}      & VARCHAR(100) & NN     & Acción realizada. \\
\texttt{detalle}     & TEXT         & ---    & Detalle del evento. \\
\texttt{modulo}      & VARCHAR(50)  & ---    & Módulo del sistema. \\
\texttt{entidad}     & VARCHAR(50)  & ---    & Entidad afectada. \\
\texttt{equipoCodigo}& VARCHAR(30)  & ---    & Código de equipo relacionado. \\
\texttt{tipoAccion}  & VARCHAR(20)  & ---    & Tipo de acción (Crear, Actualizar, etc.). \\
\texttt{estado}      & VARCHAR(20)  & NN, DF: 'Éxito' & Resultado del evento. \\
\texttt{ipOrigen}    & VARCHAR(45)  & ---    & Dirección IP de origen. \\
\bottomrule
\end{longtable}

\textbf{Relaciones:}
\begin{itemize}[nosep]
  \item N $\rightarrow$ 1 con \texttt{Usuario} (\texttt{usuario}).
\end{itemize}
\end{tabladesc}

%============================================================
% TABLA 24
%============================================================
\begin{tabladesc}{Auditoría de Cambios}{audit\_changes}{Auditoría}

Registro detallado de cambios en campos específicos de las entidades.

\bigskip
\begin{longtable}{L{3.5cm} L{2.5cm} L{2cm} L{7.5cm}}
\toprule
\textbf{Atributo} & \textbf{Tipo} & \textbf{Restricciones} & \textbf{Descripción} \\
\midrule
\endhead
\bottomrule
\endfoot
\texttt{id}          & VARCHAR(30)  & PK, NN & Identificador único. \\
\texttt{tabla}       & VARCHAR(50)  & NN     & Nombre de la tabla modificada. \\
\texttt{operacion}   & VARCHAR(10)  & NN     & Tipo de operación (INSERT, UPDATE, DELETE). \\
\texttt{registroId}  & VARCHAR(50)  & ---    & Identificador del registro afectado. \\
\texttt{usuarioId}   & VARCHAR(20)  & FK     & Referencia a \texttt{Usuario.id}. \\
\texttt{campo}       & VARCHAR(50)  & ---    & Nombre del campo modificado. \\
\texttt{valorAnterior}& TEXT        & ---    & Valor previo al cambio. \\
\texttt{valorNuevo}  & TEXT         & ---    & Valor posterior al cambio. \\
\texttt{timestamp}   & TIMESTAMP    & NN, DF: now & Momento del cambio. \\
\bottomrule
\end{longtable}

\textbf{Relaciones:}
\begin{itemize}[nosep]
  \item N $\rightarrow$ 1 con \texttt{Usuario} (\texttt{usuario}).
\end{itemize}
\end{tabladesc}

%============================================================
\newpage
\section{Diagrama Resumen de Relaciones}
%============================================================

A continuación se listan las cardinalidades entre las entidades principales:

\bigskip
\begin{longtable}{L{5cm} L{3cm} L{5cm}}
\toprule
\textbf{Entidad origen} & \textbf{Cardinalidad} & \textbf{Entidad destino} \\
\midrule
\endhead
\bottomrule
\endfoot
Edificio               & 1 $\rightarrow$ N & Laboratorio \\
Laboratorio            & N $\rightarrow$ 1 & Edificio \\
Laboratorio            & N $\rightarrow$ 1 & Persona (encargado) \\
Laboratorio            & 1 $\rightarrow$ N & Equipo \\
Laboratorio            & 1 $\rightarrow$ N & Periférico \\
Laboratorio            & 1 $\rightarrow$ N & InventarioItem \\
Laboratorio            & 1 $\rightarrow$ N & Incidencia \\
Laboratorio            & 1 $\rightarrow$ N & Mantenimiento \\
Laboratorio            & 1 $\rightarrow$ N & Asignación \\
Laboratorio            & 1 $\rightarrow$ N & ReportePasante \\
Equipo                 & N $\rightarrow$ 1 & Laboratorio \\
Equipo                 & N $\rightarrow$ 1 & EstadoEquipo \\
Equipo                 & 1 $\rightarrow$ N & Periférico \\
Equipo                 & 1 $\rightarrow$ N & Incidencia \\
Equipo                 & 1 $\rightarrow$ N & Mantenimiento \\
Equipo                 & 1 $\rightarrow$ N & InventarioItem \\
Equipo                 & 1 $\rightarrow$ N & Asignación \\
Persona                & 1 $\leftrightarrow$ 1 & Usuario \\
Persona                & 1 $\rightarrow$ N & Laboratorio (encargado) \\
Persona                & 1 $\rightarrow$ N & Incidencia \\
Usuario                & N $\rightarrow$ 1 & Rol \\
Usuario                & 1 $\rightarrow$ N & Mantenimiento (técnico) \\
Usuario                & 1 $\rightarrow$ N & Asignación (técnico) \\
Usuario                & 1 $\rightarrow$ N & ReportePasante (pasante) \\
Usuario                & 1 $\rightarrow$ N & ReportePasante (atendió) \\
Usuario                & 1 $\rightarrow$ N & Incidencia \\
Usuario                & 1 $\rightarrow$ N & Log \\
Usuario                & 1 $\rightarrow$ N & AuditChange \\
Usuario                & 1 $\rightarrow$ N & PasanteHistorial \\
Mantenimiento          & 1 $\leftrightarrow$ 1 & MantenimientoDetalle \\
MantenimientoDetalle   & 1 $\rightarrow$ N & Checklist \\
MantenimientoDetalle   & 1 $\rightarrow$ N & InsumoUsado \\
Insumo                 & 1 $\rightarrow$ N & InsumoUsado \\
CategoriaInventario    & 1 $\rightarrow$ N & InventarioItem \\
Rol                    & 1 $\rightarrow$ N & Usuario \\
TipoMantenimiento      & 1 $\rightarrow$ N & Mantenimiento \\
EstadoMantenimiento    & 1 $\rightarrow$ N & Mantenimiento \\
EstadoEquipo           & 1 $\rightarrow$ N & Equipo \\
EstadoIncidencia       & 1 $\rightarrow$ N & Incidencia \\
\bottomrule
\end{longtable}

%============================================================
\newpage
\section{Resumen de Convenciones}
%============================================================

\begin{itemize}[nosep]
  \item \textbf{IDs:} Prefijo auto-generado con sufijo numérico (\texttt{EQ-001}, \texttt{LAB-001}, etc.).
  \item \textbf{Auditoría:} Toda tabla operacional incluye \texttt{createdAt} y \texttt{updatedAt}.
  \item \textbf{Borrado lógico:} Todas las tablas operacionales usan \texttt{activo} (BOOLEAN) en lugar de DELETE físico.
  \item \textbf{Códigos de inventario:} Equipos, periféricos e ítems de inventario comparten la misma estructura de códigos (\texttt{codigoItic}, \texttt{codigoFacultativo}, \texttt{codigoUmsa}).
  \item \textbf{Cascade:} Solo en relaciones críticas: Persona $\rightarrow$ Usuario (delete), MantenimientoDetalle $\rightarrow$ Checklist/InsumoUsado (delete).
\end{itemize}

\end{document}
